\scanpage{6}

% archive sheet number: 5

On suppose donc \deleted{donné un hom de \uline{$A$-algèbres} \unsure{unitaires}} que \deleted{$f:D\longrightarrow C$ et un \uline{hom de $A$-algèbres unitaires}}

\[
\deleted{g_C:B\longrightarrow C\otimes_{\deleted{\unsure{A|D}}}C}
\]

\deleted{rendant commutatif le diagramme suivant}

\diagram{\begin{tikzcd}[ampersand replacement=\&]
{} \& \deleted{B\otimes C}\arrow[l] \& \deleted{C}\arrow[l,"{\deleted{q_C}}"']
\end{tikzcd}}

\deleted{ayant la propriété suivante :} l'hom composé

\diagram{\begin{tikzcd}[ampersand replacement=\&]
B\otimes C \& B\otimes C\otimes C\arrow[l,"{\mathrm{id}_B\otimes\Delta_C}"'] \& C\otimes C\arrow[l,"{q_C\otimes\mathrm{id}_C}"']
\end{tikzcd}}

[i.e. l'hom $C\otimes C\xrightarrow{(q_C,p_2)}B\otimes\deleted{\unsure{A}}C$] est un \uline{isomorphisme} (de $A$-algèbres). Il revient au \unsure{même} de dire qu'il \uline{existe un hom de $A$-algèbres unitaires} $\boxed{g_C:B\longrightarrow C\otimes C}$

tel que l'on ait commutativité des

\diagram{\begin{tikzcd}[ampersand replacement=\&]
(*)\quad C\otimes C \& \deleted{\unsure{B\otimes C\otimes C}}\arrow[l,"{(g_C,\deleted{\unsure{1}}p_2\deleted{\unsure{C}})}"'] \& B\otimes C\arrow[l,"{\deleted{\unsure{g\otimes p_2}}}"',"{\deleted{\unsure{p_2}}}"] \& C\arrow[l,"q_C"']\arrow[lll,bend left=17,"p_1"]
\end{tikzcd}}

\diagram{\begin{tikzcd}[ampersand replacement=\&]
(**)\quad B\otimes C\& C\otimes C\arrow[l,"{(q_C,p_2)}"']\& B\arrow[l,"g_C"']\arrow[ll,bend left=24,"p_1"]
\end{tikzcd}}

Nous allons \deleted{\unsure{voir}} \inserted{examiner} si sous ces conditions, il est bien exact que $C$ soit fibré principal \uuline{hom} sous $G$, de \deleted{\unsure{g}} base l'objet unité $A$ de $\unsure{\mathcal{C}|\mathscr{C}}$.

Supposons donc qu'on a un hom de $A$-algèbres \inserted{commutatives unitaires}

\[
h:D\longrightarrow C
\]

rendant commutatif les diagrammes
